\section{实验结果与改进路径}

\subsection{整体改进轨迹}

LangChain 的 Harness Engineering 改进过程可以总结为以下演进路径：

\begin{table}[H]
\centering
\caption{Harness 迭代改进路径}
\begin{tabular}{p{1.5cm}p{6cm}p{2cm}p{2.5cm}}
\toprule
\textbf{阶段} & \textbf{改进内容} & \textbf{得分} & \textbf{增幅} \\
\midrule
Baseline & 默认 Prompt + 标准 Tools + 标准 Middleware & 52.8\% & --- \\
+SV & 加入 Self-Verification 四阶段框架 + PreCompletionChecklist & --- & 显著 \\
+CE & 加入 Context Engineering（目录、工具、时间预算） & --- & 中等 \\
+LD & 加入 LoopDetectionMiddleware & --- & 中等 \\
+RS & 采用 Reasoning Sandwich 策略 & 66.5\% & --- \\
\midrule
\multicolumn{2}{l}{\textbf{总提升}} & \multicolumn{2}{l}{\textbf{+13.7 百分点}} \\
\bottomrule
\end{tabular}
\end{table}

值得注意的是，这 13.7 个百分点的提升\textbf{完全来自 Harness 改进}，模型（GPT-5.2-Codex）保持不变。

\subsection{跨模型适配性}

LangChain 还用早期版本的 Harness 测试了 Claude Opus 4.6，得分为 59.6\%。这个分数具有竞争力，但低于针对 Codex 优化后的 Harness 的表现。原因并非 Claude 模型本身不行，而是\textbf{没有为 Claude 运行同样的 Harness 改进循环}。

\begin{warningbox}{Harness 并非模型通用}
不同模型需要不同的 Prompting 策略（参见 Codex 和 Claude 各自的 Prompting Guide）。虽然某些原则（如重视验证、良好的上下文准备）具有通用性，但\textbf{针对每个模型跑几轮 Harness 迭代是必要的}。一个为 Model A 优化的 Harness 直接用于 Model B，通常不会达到最优表现。
\end{warningbox}

\subsection{本章小结}

实验结果清晰地展示了 Harness Engineering 的投资回报：仅改动外部系统就能带来 13.7 个百分点的提升。但 Harness 改进并非一劳永逸——它需要针对特定模型做迭代优化，不同模型需要不同的 Harness 配置。


